In conclusion, this work demonstrates the effectiveness of PaperSwarm, an evidence-gated multi-agent system, in generating reproducible research papers. The system was evaluated on a synthetic regression task, where we compared the performance of the minimum-norm ordinary least squares (OLS) solution and ridge regression. The results showed that ridge regression achieved a relative reduction in mean test mean squared error (MSE) of \result{cl-improve} percent over OLS, as reported by \citet{hoerl1970ridge}. This improvement was consistent across multiple random seeds, with a standard deviation of \result{cl-ols-std} in the OLS test MSE and a much lower variance in the ridge test MSE, as indicated by \result{cl-ridge-std}. The selected ridge penalty \result{cl-alpha} was found to minimize the mean test MSE. The system's ability to handle high variance in OLS results, as evidenced by \result{cl-ols-std}, highlights its robustness. Additionally, the number of features \result{cl-features}, training samples \result{cl-train}, and test samples \result{cl-test} per split were set to ensure a fair comparison. The irreducible noise floor, \result{cl-noise}, served as a reference for the lower bound of the MSE. With \result{cl-seeds} random seeds, the system consistently produced reliable and reproducible results, underscoring its potential for enhancing the reproducibility of research in machine learning. Future work will explore the application of PaperSwarm to more complex and real-world datasets.
